\chapter{Introduction}
\label{ch:introduction}

This is a placeholder chapter for the AutoTex book generation system.

\section{Background}
\label{sec:background}

Background content will be generated here. The AutoTex system uses AI agents to produce
high-quality scientific content in \LaTeX{}, including mathematical formulas, tables,
diagrams, and illustrations.

\begin{definitionbox}[Key Term]
  This is an example of a definition box provided by the \texttt{autotex-book} style.
  Definition boxes are used to highlight important terms and concepts throughout the book.
\end{definitionbox}

\section{Objectives}
\label{sec:objectives}

The main objectives of this book are:

\begin{enumerate}
  \item First objective
  \item Second objective
  \item Third objective
\end{enumerate}

\begin{theorembox}[Sample Theorem]
  For all $n \in \mathbb{N}$, the sum of the first $n$ positive integers is given by:
  \[
    \sum_{k=1}^{n} k = \frac{n(n+1)}{2}
  \]
\end{theorembox}

\section{Overview}
\label{sec:overview}

\begin{examplebox}[Example Usage]
  This box demonstrates the example environment. It can be used to present
  worked examples, case studies, or illustrative scenarios.
\end{examplebox}

\begin{notebox}[Important Note]
  This is a note box for warnings, important remarks, or caveats that
  the reader should pay attention to.
\end{notebox}

% Example table using booktabs
\begin{table}[htbp]
  \centering
  \begin{tabular}{lcc}
    \toprule
    Feature & Status & Priority \\
    \midrule
    Text generation & Active & High \\
    Image generation & Active & High \\
    PDF review & Active & Medium \\
    \bottomrule
  \end{tabular}
  \caption{Example table demonstrating booktabs formatting.}
  \label{tab:example}
\end{table}

% Example figure inclusion (uncomment when image exists)
% \begin{figure}[htbp]
%   \centering
%   \includegraphics[width=0.8\textwidth]{example-figure}
%   \caption{An example figure.}
%   \label{fig:example}
% \end{figure}
